\section{Calibration and randomized opposed launches}
\label{sec:calibration}
The scale-independent estimate in Proposition~\ref{prop:mean-stability}
concerns the propagation of mean error, not the implementation of an
input law. We first bound physical preparation errors without excluding
grazing commands. We then allow a calibrated distribution of directions
and durations rather than a single collimated displacement.

\subsection{A geometric bound for imperfect preparation}
For a compact convex body $C$ let
$S_a(C)=C+[-a,0]$. A segment $[q,q+a]$ hits $C$ precisely when
$q\in S_a(C)$. For a locally finite union of bodies this gives
\begin{equation}\label{eq:morph-bit}
 B_a(q)=\one_{\cup_C S_a(C)}(q)-\one_{\mathcal O}(q).
\end{equation}
The subtraction is essential: the first indicator counts a solid start
as a hit, whereas our recorded bit counts it as zero.

\begin{lemma}[Local convex boundary tubes]\label{lem:tube}
There is an absolute constant $c_0$ such that, for every planar compact
convex body $C$, disk $B_s(x)$ and $r>0$,
\[
 \area\{q\in B_s(x):\dist(q,\partial C)\le r\}
                 \le c_0(s+r)r.
\]
The bound also applies to the convex swept body $S_a(C)$, irrespective
of its straight sides or nonsmooth junctions.
\end{lemma}
\begin{proof}
Use signed distance, negative inside $C$, and integrate its level sets
between $-r$ and $r$ by the coarea formula. Its gradient has length one
almost everywhere away from the boundary and the interior medial set.
Positive levels are boundaries of outer parallel bodies; negative
levels are boundaries of convex erosions, until the erosion becomes
empty. The part of any such convex boundary in $B_s(x)$ has length at
most $2\pi s$: it is part of the boundary of its intersection with the
disk, whose perimeter is at most that of the disk by the planar Cauchy
perimeter formula. Coincident boundary pieces only reduce this bound.
Approximation by slightly larger disks handles tangencies and exceptional
levels. Integration gives $4\pi sr$; the displayed weaker bound also
covers all degenerate limiting cases.
\end{proof}

\begin{theorem}[Position, duration and angle calibration]\label{thm:calibration}
Assume the planar components are convex and mutually separated by
$d_0>0$. A nominal command has density $f_{x,\sigma}$, displacement
$a=tv$, $|v|=1$, $t\le b$, and $0<\sigma\le\sigma_0$.
Suppose an actual attempt can be coupled to this command so that
\[
 |q'-q|\le\ell,\qquad |t'-t|\le\tau,\qquad
 \angle(v',v)\le\alpha,
\]
and remains a unit-speed straight flight until its first collision.
Both nominal and actual protocols assign zero to solid starts and free
misses and count all attempts. Set
\begin{equation}\label{eq:calibration-size}
             r=\ell+\tau+b\alpha.
\end{equation}
If $r\le\sigma$, then
\begin{equation}\label{eq:calibration-bias}
 \left|\E B_{t'v'}(q')-P_a(f_{x,\sigma})\right|
                  \le C_{d_0,b,\sigma_0,\kappa}\,r/\sigma.
\end{equation}
The coupling errors may depend on $q$ and need not have mean zero.
An exceptional coupling event of probability $\zeta$ adds at most
$\zeta$ to the bound. The estimate is uniform in the nominal direction
and includes commands tangent to obstacles.
\end{theorem}
\begin{proof}
One has $|t'v'-tv|\le\tau+b\alpha$. The Hausdorff distance between
$S_{t'v'}(C)$ and $S_a(C)$ is at most this quantity. For convex bodies
$K,K'$ with $d_H(K,K')\le s$, the inclusions
$K_{-s}\subset K'\subset K+s\overline B$ follow directly from their
support inequalities. Thus a change of swept-body membership, including
the start displacement, can occur only when the nominal $q$ is within
$r$ of $\partial S_a(C)$. A change of solid-start membership is similarly
confined to the $r$-tube of $\partial C$. These containments hold for
every allowed actual displacement, even if its error depends on $q$.

Only a uniformly bounded number $N_0=N_0(d_0,b,\sigma_0)$ of components
can contribute such tubes in $B_\sigma(x)$. Indeed each contributing
component has a point in $B_{\sigma+b+r}(x)$; chosen points from distinct
components are separated by $d_0$, so disjoint radius-$d_0/2$ disks give
a uniform packing bound. Apply Lemma~\ref{lem:tube} to these bodies
and their sweeps in \eqref{eq:morph-bit}. Their exceptional area is at
most $C N_0(\sigma+r)r$. Since
$\|f_{x,\sigma}\|_\infty=\sigma^{-2}\|\kappa\|_\infty$ and
$r\le\sigma$, its probability is at most $Cr/\sigma$. Outside it
both bits agree. A failed coupling contributes at most its probability
because the bits belong to $[0,1]$.
\end{proof}

A systematic center error, angular error and timing error may therefore
all be charged explicitly to the mean-error budget. For fixed allowed
mean bias $b_0$, the sufficient condition is
$\ell+\tau+b\alpha\le c b_0\sigma$, not a scale-free apparatus
tolerance. No calibration estimate is obtained from the collision bits
themselves. Arbitrary curved pre-impact paths or an unknown distribution
outside this coupling model are not included.

\subsection{From one displacement to a stopped averaging operator}
Let $\rho$ be a known probability measure on displacements. In the
forward protocol draw $a\sim\rho$ and $q\sim f$ independently and
attempt $(q,a)$. In the reverse protocol draw the same law for $a$, draw
$q\sim f$, and attempt $(q+a,-a)$. These two protocols need not use
the same physical particle or the same random draw. Their prescribed
joint distributions must, however, have this translated relationship.
Only their mean bits are used. Write their means as $P_\rho^+(f)$ and
$P_\rho^-(f)$, and define
\[
 (T_\rho u)(x)=\int u(x+a)\rho(\dd a).
\]
Averaging Theorem~\ref{thm:balance} shows that their difference is
$\int f(T_\rho\chi-\chi)$; for the kernel commands it is
\begin{equation}\label{eq:random-balance}
       g_\sigma(x)=(T_\rho-I)u_\sigma(x).
\end{equation}
This is a coupled preparation design. Independently randomized reverse
starts without the translation by their sampled displacement do not
in general give \eqref{eq:random-balance}.

\begin{theorem}[Finite stopped inverse]\label{thm:stopped}
Let $0\le u\le1$, let $T$ be the transition operator of a Markov chain
$(X_j)$, and suppose that, for every starting point in the region in
question, the first visit to $\{u=0\}$ occurs by time $m$ almost surely.
Values of $u$ and $g=(T-I)u$ need only be defined on its $m$-step
reachable region. Put
\begin{equation}\label{eq:bellman}
 V_0=0,\qquad V_{n+1}=\max\{0,TV_n-g\},\qquad 0\le n<m.
\end{equation}
Then $V_m=u$. For a perturbation $\widehat g$ with uniform error
at most $\xi$, the same recursion, optionally clipped to $[0,1]$,
has error at most $m\xi$ at step $m$. Each additional uniform error
in evaluating $TV_n$ adds at most its magnitude to the final bound.

For $T=T_\rho$, the hypothesis holds for $u=\chi$ and $u=u_\sigma$
under the component bounds if, for some unit $v$ and $\gamma>0$,
\begin{equation}\label{eq:cone}
 a\cdot v\ge\gamma,
 \quad |a|\le b<d_0-2\sigma_0
 \quad(\rho\text{-almost every }a),
 \qquad m\gamma>D_0+2\sigma_0.
\end{equation}
Thus two averaged scalar-launch functionals with a fixed, possibly
continuous spread of durations and directions determine the local
obstacle set. Neither a realized impact position nor a free reference
point is supplied. The aperture enlargement is at most $mb+\sigma_0$.
\end{theorem}
\begin{proof}
Induction gives $0\le V_n\le u$, because $Tu-g=u$. At a point where
$u=0$, every $V_n$ is zero. At a point where $u>0$, writing
$e_n=u-V_n\ge0$ gives
$e_{n+1}=\min(u,Te_n)\le Te_n$. Consequently
\[
 e_m(x)\le\E_x\left[u(X_m)
                 \prod_{j=0}^{m-1}\one_{\{u(X_j)>0\}}\right]=0.
\]
The last equality uses the assumed visit by time $m$, including its
endpoint. Positivity of $T$ and nonexpansiveness of the maximum prove
$\|V_{n+1}-\widehat V_{n+1}\|_\infty
 \le\|V_n-\widehat V_n\|_\infty+\xi$.
Clipping cannot increase error relative to $V_n\in[0,1]$. This proves
the stability assertions.

Under \eqref{eq:cone}, a path whose first $m+1$ points have positive
$u_\sigma$ must stay in a single $\sigma$-enlarged component: successive
steps are shorter than the distance between distinct enlarged components.
Its endpoints differ by at least $m\gamma$ in direction $v$, exceeding
the diameter of that component. This is impossible. The indicator case
uses the same argument without enlargement. Approximate-identity
convergence after \eqref{eq:random-balance} recovers $\chi$.
\end{proof}

The recursion is the elementary finite-horizon obstacle recursion, here
with a physically supplied forcing $g$ and a geometrically bounded
zero-hitting time. For $\rho=\delta_a$ its expansion is exactly the
minimum of partial sums in \eqref{eq:min-inverse}. The finite stopping
condition cannot be discarded: for a symmetric nearest-neighbor walk
and $u=1$ on two adjacent sites and zero elsewhere, the two nonzero
values of $V_n$ are $1-2^{-n}$, not one at any finite $n$.
Theorem~\ref{thm:calibration} may be conditioned on the sampled nominal
displacement in either protocol. Uniformly bounded implementation errors
then perturb each averaged scalar mean by the same explicit bound.
